% --------------------------------------------------------------------------
% DMRN one-page paper: SCRAPL in JAX
% Built from template.tex; uses dmrn2025.sty, IEEEtran.bst, scrapl_jax_dmrn.bib
% Placeholders are marked with \tbd (rendered in red).
% --------------------------------------------------------------------------

\documentclass{article}
\usepackage{caption}
\usepackage{times,amsmath,amssymb,graphicx,url,dmrn2025}
\usepackage{indentfirst}
\usepackage{booktabs,multirow}
\usepackage{xcolor}
\setlength{\parskip}{0.5em}

% smaller font size for captions
\captionsetup[figure]{font=footnotesize,labelfont=footnotesize}
\captionsetup[table]{font=footnotesize,labelfont=footnotesize}

% reduce blank space after figures
\setlength{\belowcaptionskip}{-10pt}

% roman numbering of sections
\renewcommand{\thesection}{\Roman{section}}

% Placeholder for values and claims still to be filled in / verified
\newcommand{\tbd}[1][TBD]{\textcolor{red}{#1}}

% Four authors, two affiliations (same layout as \twoauthors in dmrn2025.sty)
% #1-#4: author names, #5: affiliation 1, #6: email, #7: affiliation 2
\makeatletter
\def\fourauthorstwoaffils#1#2#3#4#5#6#7{\gdef\@address{}
   \gdef\@name{{\underline{#1}\textsuperscript{1}, #2\textsuperscript{2},
     #3\textsuperscript{1} and #4\textsuperscript{2}}\\\\
   \normalsize{\textsuperscript{1}#5, \underline{#6}}\\\normalsize{\textsuperscript{2}#7}}}
\makeatother

% Title.
% --------------------
\title{SCRAPL in JAX: Efficient Scattering Transform Losses}

% Single addresses (uncomment and modify for single-address case).
% --------------------
\oneauthor
{Christopher Mitcheltree}
{Neutone Inc., Tokyo, Japan}
{christopher@neutone.jp}

\begin{document}
\tenpt
\maketitle

\begin{sloppy}
% Pull the abstract up so its first line is flush with the Table 1 caption,
% which \belowcaptionskip{-10pt} lifts above the column top. Tuned by measuring
% the rendered PDF; re-tune if the caption skips or abstract spacing change.
\vspace*{-32.9pt}
\begin{abstract}
``\underline{Sc}attering transform with \underline{Ra}ndom \underline{P}aths for machine \underline{L}earning'' (SCRAPL) is a stochastic optimization scheme that makes wavelet scattering transforms tractable as differentiable loss functions for perceptual quality assessment of deep inverse problems in computer vision, speech, and audio processing.
We present a JAX implementation of SCRAPL for the joint time--frequency scattering transform (JTFS).
On GPU, our implementation is ${\sim}10\times$ faster than PyTorch while using less memory.
\end{abstract}

%for perceptual quality assessment of deep inverse problems in computer vision, speech, and audio processing.

% The DMRN example papers omit keywords, so the keywords environment is left out.

% ==========================================================================
\vspace{-8pt}
\section{SCRAPL}
\label{sec:scrapl}
\vspace{-5pt}

A wavelet scattering transform (ST) decomposes a signal $\boldsymbol{x}$ into $P$ coefficient vectors $\phi_p(\boldsymbol{x})$, known as \emph{paths}.
The Euclidean distance between ST coefficients provides informative gradients for inverse problems in vision, speech, and audio, but evaluating all $P$ paths makes it an expensive training loss.
SCRAPL~\cite{mitcheltree2026scrapl} instead samples one path per training step.
For a reconstruction $\tilde{\boldsymbol{x}} = F_{\boldsymbol{w}}(\boldsymbol{x})$ by a model with weights $\boldsymbol{w}$, the ST loss is an expectation over paths:
\vspace{-5pt}
\begin{equation}
\mathcal{L} = \sum_{p=0}^{P-1}
  \big\Vert \phi_p(\boldsymbol{x}) - \phi_p(\tilde{\boldsymbol{x}}) \big\Vert_2^2
  = \mathbb{E}_{p\sim\mathcal{U}_P}\big[\mathcal{L}_p\big],
\label{eq:scrapl}
\vspace{-5pt}
\end{equation}
where $\mathcal{L}_p = P\,\Vert \phi_p(\boldsymbol{x}) - \phi_p(\tilde{\boldsymbol{x}}) \Vert_2^2$ and $\mathcal{U}_P$ is uniform over paths.
Hence, $\nabla_{\boldsymbol{w}}\mathcal{L}_p$ is an unbiased estimate of $\nabla_{\boldsymbol{w}}\mathcal{L}$ at roughly $1/P$ of the cost.
SCRAPL reduces the variance of this estimate with $\mathcal{P}$-Adam, which keeps separate Adam moments per path; $\mathcal{P}$-SAGA, a path-wise version of SAGA~\cite{defazio2014saga}; and $\theta$-importance sampling ($\theta$-IS), which weights each path by the curvature of its loss landscape, estimated from its sensitivity to the parameters $\boldsymbol{\theta}$ of a signal model, such as a synthesizer.
As in~\cite{mitcheltree2026scrapl}, we implement SCRAPL for the joint time--frequency scattering transform (JTFS)~\cite{anden2019joint},
%which demodulates spectrotemporal patterns at multiple scales and rates and
which is an informative perceptual audio loss for differentiable digital signal processing (DDSP).

% ==========================================================================
\vspace{-8pt}
\section{JAX Implementation}
\label{sec:jax}
\vspace{-5pt}

The \texttt{scrapl-loss} Python package\footnote{\url{https://github.com/christhetree/scrapl}} now includes a JAX backend that does not depend on PyTorch.
%, complementing JAX audio tools such as Korvax~\cite{marttila2025korvax}.

\noindent \textbf{Single-path JTFS:} Each path is computed by a single-path JTFS built on Kymatio's JAX backend~\cite{andreux2020kymatio}.
Paths differ in their filters and output shapes, so each path is just-in-time (JIT) compiled separately with a static path index.

\noindent \textbf{Stateless loss:} \texttt{SCRAPLLoss} is an immutable Equinox module: paths are sampled from a JAX PRNG key, and all training state lives in caller-managed PyTrees.

\noindent \textbf{$\mathcal{P}$-Adam and $\mathcal{P}$-SAGA:} Both are Optax transformations with per-path state, chainable with standard Optax transforms like weight decay and learning-rate scaling.

\noindent \textbf{$\theta$-IS:} Improving on PyTorch, Hessian-vector products are computed forward-over-reverse (\texttt{jax.jvp} of \texttt{jax.vjp}), power iteration runs in \texttt{lax.while\_loop} with early stopping, batches are accumulated with \texttt{lax.scan}, and synthesizer parameters can be processed in parallel with \texttt{vmap}.
Path PRNG keys are derived from the path index, so the warmup can be split across devices with identical results.

% ==========================================================================
\vspace{-8pt}
\section{Benchmarks}
\label{sec:bench}
\vspace{-5pt}


% For reference, Table 5 of the ICLR paper (PyTorch 2.8.0, 1 NVIDIA RTX A5000,
% CUDA 12.4, forward + backward, same signal and JTFS settings):
%   JTFS:   1731 ms, 12967 MB
%   SCRAPL: 89.82 ms (median over paths), 2503 MB (max over paths)
% Compile times in the text are from benchmarks/min_run_1_sec (compile_time column:
% scrapl_mean, warmup scrapl_mean, jtfs of the JAX TSVs).
% Numbers from benchmarks/min_run_15_sec/*_cuda.tsv: JTFS = "jtfs" row of
% benchmark_{jax,torch}_cuda.tsv, SCRAPL = "scrapl_mean" row of the same files,
% theta-IS = "scrapl_mean" row of benchmark_warmup_{jax,torch}_cuda.tsv.
% FLOPs are the JAX (XLA) counts; none are measured for the theta-IS warmup.
\begin{table}[t]
  \centering
  \caption{Benchmark results for one optimization step (forward + backward, 44100 samples of audio, batch size 1, 1 thread, single precision, NVIDIA RTX A5000 GPU, CUDA 12.8, PyTorch 2.10, JAX 0.12).
  For $\theta$-IS, one step estimates the curvature for one path (10 power iterations, $\dim(\boldsymbol{\theta}) = 1$).
  SCRAPL and $\theta$-IS are benchmarked for $P=315$ paths and averaged, except for memory (maximum).
  JAX times exclude JIT compilation; FLOPs are counted by XLA and not measured for $\theta$-IS.}
  \vspace{-8pt}
  \label{tab:bench}
  \footnotesize
  \setlength{\tabcolsep}{3.5pt}
  \begin{tabular}{llcccc}
    \toprule
    \textbf{Method} & \textbf{Framework} & \textbf{FLOPs} $\downarrow$
      & \textbf{Time} $\downarrow$ & \textbf{IQR} $\downarrow$
      & \textbf{Mem.} $\downarrow$ \\
    & & & \textbf{(ms)} & \textbf{(ms)} & \textbf{(MB)} \\
    \midrule
    \multirow{2}{*}{\begin{tabular}[c]{@{}l@{}}Full-tree\\(JTFS)\end{tabular}}
      & PyTorch & \multirow{2}{*}{$2.91 \times 10^{11}$}
                          &     1991 &     257 &     6164 \\
      & JAX     &         & \bf  328 & \bf 27.1 & \bf 5562 \\
    \midrule
    \multirow{2}{*}{SCRAPL}
      & PyTorch & \multirow{2}{*}{$3.26 \times 10^{9}$}
                          &      136 &     24.5 &      680 \\
      & JAX     &         & \bf 9.77 & \bf 2.03 & \bf  314 \\
    \midrule
    \multirow{2}{*}{$\theta$-IS}
      & PyTorch & \multirow{2}{*}{--}
                          &      968 &      120 &      730 \\
      & JAX     &         & \bf  103 & \bf 6.78 & \bf  638 \\
    \bottomrule
  \end{tabular}
  \vspace{-8pt}
\end{table}

Table~\ref{tab:bench} compares both implementations using the JTFS hyperparameters of~\cite{mitcheltree2026scrapl}.
JAX is 6.1$\times$, 13.9$\times$, and 9.4$\times$ faster than PyTorch for the full-tree JTFS, SCRAPL, and $\theta$-IS, respectively, and uses less memory throughout.
With JAX, a SCRAPL step is 34$\times$ faster than a full-tree JTFS step.

XLA compiles each path when it is first sampled, taking 7.8\,s per path on average (11\,s for $\theta$-IS) versus 119\,s for the full-tree JTFS.
This is paid once per path and can be reused across runs with JAX's persistent compilation cache.
We omit \texttt{torch.compile}, as it compiled more than 10$\times$ slower than JAX.

% -------------------------------------------------------------------------
% The bibliography style prints its own numbered "References" section.
\bibliographystyle{IEEEtran}
\vspace{-5pt}
{\footnotesize\bibliography{scrapl_jax_dmrn}}

\end{sloppy}
\end{document}
